% DeepINSIGHT report, v6.
% Build:  python -m deepinsight paper   (or: cd paper && tectonic main.tex)
% Every number is traced to a recorded output file in NUMBERS.md. File-level detail (paths, step
% codes, file formats) lives in docs/DATA_FLOW.md, not here. Notes for the authors are
% \authornote{...} (red); \renewcommand{\authornote}[1]{} hides them.
\documentclass[11pt]{article}
\usepackage[a4paper,margin=1.9cm]{geometry}
\usepackage[T1]{fontenc}
\usepackage{lmodern}
\usepackage{microtype}
\usepackage{amsmath,amssymb}
\usepackage{graphicx}
\usepackage{booktabs,tabularx}
\usepackage{multicol}
\usepackage{xcolor}
\usepackage[numbers,sort&compress]{natbib}
\usepackage{caption}
\usepackage[compact]{titlesec}
\usepackage[colorlinks=true,linkcolor=black,citecolor=black,urlcolor=blue!50!black]{hyperref}
\captionsetup{font=small,labelfont=bf,labelsep=period}
\makeatletter\g@addto@macro\UrlBreaks{\do\_\do\-\do\.}\makeatother
\newcommand{\file}[1]{\nolinkurl{#1}}
\newcommand{\authornote}[1]{{\color{red!70!black}\small [Authors: #1]}}
\setlength{\textfloatsep}{10pt plus 2pt minus 2pt}
\setlength{\parskip}{2pt}
\renewcommand{\topfraction}{0.9}\renewcommand{\bottomfraction}{0.8}\renewcommand{\textfraction}{0.1}\renewcommand{\floatpagefraction}{0.8}
\setcounter{topnumber}{3}\setcounter{totalnumber}{4}
\titlespacing*{\paragraph}{0pt}{3pt}{0.6em}

\begin{document}

\begin{center}
{\LARGE\bfseries DeepINSIGHT: A Temporal Regulatory Disentanglement Framework\\[2pt] for Non-Coding Variant Interpretation\par}
\vspace{8pt}
{\large Julien V. Mutton$^{1}$, Troy M. LaPolice$^{2}$, Anthony Kwong$^{1}$, Shaun Mahoney$^{3}$, Christian D. Huber$^{1,*}$\par}
\vspace{5pt}
{\small
$^{1}$Pennsylvania State University, Department of Biology, University Park, PA 16802, USA\\
$^{2}$Princeton University, Department of Ecology \& Evolutionary Biology, Princeton, NJ 08544, USA\\
$^{3}$Pennsylvania State University, Department of Biochemistry and Molecular Biology, University Park, PA 16802, USA\\[3pt]
$^{*}$Correspondence: Christian D. Huber, M\"uller Laboratory, Room 512A, University Park, PA 16802, USA.
Email: \href{mailto:cdh5313@psu.edu}{cdh5313@psu.edu}\par}
\end{center}

\authornote{Title: no analysis supports ``Temporal'' or ``Disentanglement''. All results use a 0.1\% sample of sites.}

%==================================================================================================
\section*{Abstract}
Most variants linked to human disease lie outside genes, where their effects are hard to predict.
LINSIGHT scores each non-coding site by the probability that mutations there are removed by natural
selection, learned from human genetic variation and human--chimpanzee differences through 48
pre-computed genomic annotations, without disease labels. Genomic language models read DNA sequence
directly, but whether they add anything to LINSIGHT has not been tested. We rebuilt LINSIGHT from its
original data and code, then gave it one extra input: either a HyenaDNA embedding of the DNA around
each site, or a short code of that embedding learned by a sparse autoencoder. Nothing else changed.
Each extra input was compared with control inputs of the same size, such as random numbers. The
rebuilt LINSIGHT reproduced the published model. On the HGMD disease-variant benchmark, neither
sequence input improved on LINSIGHT; the full embedding made it worse. Neither input helped the model
predict evolutionary data on withheld parts of the genome. Simulations show that this second test is
too weak at our sample size to detect a modest effect, so a larger sample is needed before the
evolutionary question can be answered.

%==================================================================================================
\section*{Introduction}
Most disease-associated variants are non-coding and act on gene regulation, so their effects are
hard to read. One family of methods combines pre-computed annotations. CADD and GWAVA train classifiers on them. LINSIGHT instead fits them to evolutionary data. It builds on INSIGHT, a population-genetic model that compares variation within humans
(polymorphism) and differences from the human--chimpanzee ancestor (divergence) at a set of sites with
the same quantities at nearby neutral sites. A site under purifying selection shows fewer variants
than its neutral neighbours. LINSIGHT turns the probability that a site is under selection, written
$\rho$, into a function of 48 annotations: conservation scores, predicted binding sites and regional
functional genomics. Because it is never shown disease labels, it needs no training set of pathogenic
variants. It outperformed CADD and GWAVA on regulatory variants from HGMD, the Human Gene Mutation
Database. It never reads the DNA sequence itself, so
anything the 48 annotations miss is invisible to it.

A second family reads sequence: supervised models learn pathogenicity or experimental tracks from DNA, and genomic language models such as HyenaDNA learn the statistics of the genome without labels. These models do not use
population-genetic data, although for coding variants combining a language model with population data
has worked well.

We ask one question: \emph{does a genomic language model's reading of the DNA around a site add
information to LINSIGHT?} We keep LINSIGHT's model and training objective and change only its input. Each model gets the 48 annotations plus at most one added block: the
HyenaDNA embedding, or a sparse autoencoder code of it, or a control block of the same size. We judge
each model in two ways: on the HGMD benchmark that LINSIGHT was published on, and by how well it
predicts evolutionary data in parts of the genome held back from training.


%==================================================================================================
\section*{Results}

\subsection*{The rebuilt LINSIGHT reproduces the published model}
Our reader of the INSIGHT data gives exactly the same numbers as the authors' preparation program,
and our training code matches the authors' program parameter for parameter at every training pass on
test data. Run on all sites, the
authors' program reproduces the AUROCs of the published score track to four decimals
(Table~\ref{tab:hgmd}, second row). Retrained by us on a 0.1\% sample of sites, LINSIGHT matched the published model on the
region-matched negative set and was about 0.01 lower on the other two (Table~\ref{tab:hgmd}).

The authors' procedure runs a fixed 20 passes over the data. On our sample it stops far from the best
fit: the weights point the same way as the published ones but are much smaller (14\% of their size).
We therefore also fitted every model to convergence (``exact fit'', Methods) for the withheld-genome
test.

\subsection*{HGMD benchmark: neither sequence input improves LINSIGHT}
The benchmark asks whether a model scores 1,564 disease-causing regulatory variants from HGMD above
common variants. There are three sets of common variants (negatives): \emph{unmatched}, random
across the genome; \emph{TSS-matched}, at the same distance from a gene start as the disease variants;
and \emph{region-matched}, within 1\,kb of a disease variant. The matched sets are harder because they
remove easy location clues. We measure accuracy by the AUROC, the probability that a random disease
variant scores higher than a random negative (0.5 is chance, 1 is perfect). The benchmark was read once,
after the analysis plan was fixed (Methods).

No model with a sequence input scored above LINSIGHT on any negative set (Table~\ref{tab:hgmd}). The full embedding lowered the AUROC on all three sets, and on two of them it
did worse than 256 random numbers, so its information made the model worse than noise. The
autoencoder code lowered the AUROC on the unmatched and TSS-matched sets and was level with LINSIGHT on
the region-matched set. It did better than the full embedding on every set, but it was not better than
64 random numbers on any set, so its advantage over the embedding can be explained by its having fewer
columns. The positive control, which moves one of LINSIGHT's own annotations into the added
block, stayed within 0.004 of LINSIGHT, which shows that the added block can carry a real signal.

\begin{table}[t]
\centering\small
\caption{\textbf{HGMD AUROC of every model}, by negative set (models fitted with the authors'
procedure; score = mean $\rho$ of three training runs). Higher is better. The first two rows are the
published model; the others were trained on the 0.1\% sample. Italics: control models.}
\label{tab:hgmd}
\begin{tabular}{@{}lccc@{}}
\toprule
Model & Unmatched & TSS-matched & Region-matched \\
{\footnotesize (number of negatives)} & {\footnotesize (159,765)} & {\footnotesize (15,346)} & {\footnotesize (4,722)} \\
\midrule
Published LINSIGHT, as reported by Huang et al.\ & 0.897 & 0.759 & 0.660 \\
Published LINSIGHT score track, scored by us & 0.898 & 0.762 & 0.665 \\
\textbf{LINSIGHT (48 features), retrained} & \textbf{0.887} & \textbf{0.748} & \textbf{0.664} \\
\quad \emph{+ positive control} & 0.884 & 0.745 & 0.663 \\
\quad + HyenaDNA embedding (256) & 0.725 & 0.683 & 0.647 \\
\quad \emph{+ random columns (256)} & 0.827 & 0.719 & 0.654 \\
\quad + autoencoder code (64) & 0.824 & 0.731 & 0.665 \\
\quad \emph{+ random columns (64)} & 0.849 & 0.730 & 0.657 \\
\bottomrule
\end{tabular}
\end{table}


\subsection*{Withheld genome: no sequence input improves the fit}
This test needs no disease labels. Each model is fitted on half of the genome, split into 1-Mb blocks.
It is then asked how probable it finds the real polymorphism and divergence data on the withheld blocks.
A model that has learned something true about selection gives the withheld data a higher probability
than LINSIGHT does. No sequence input made the withheld data more probable than LINSIGHT did (Table~\ref{tab:heldout}). The
embedding did slightly worse than LINSIGHT, and significantly worse than 256 random numbers. None of the
three autoencoder codes differed from LINSIGHT or from its controls. The positive control was
indistinguishable from LINSIGHT, as it should be, because it carries the same information.

One model without any sequence input did change the fit. Replacing the 48 annotations with a 32-number
autoencoder code of the annotations themselves made the withheld data more probable than LINSIGHT
did. On the separate set of sites used for choosing models, the same model was worse than LINSIGHT.
The two results disagree, so we treat this as a result to re-test on a larger sample.

\begin{table}[t]
\centering\small
\caption{\textbf{Withheld-genome test} (exact fit). Estimate: whether the model gave the real data on
612,686 withheld sites a higher (better) or lower (worse) probability than the model it is compared
with. $p$: two-sided, paired bootstrap over 1,187 withheld 1-Mb blocks, not adjusted for the 22
comparisons. Verdict: better or worse only if $p<0.05$.}
\label{tab:heldout}\setlength{\tabcolsep}{5pt}
\begin{tabular}{@{}llccc@{}}
\toprule
Model & Compared with & Estimate & $p$ & Verdict \\
\midrule
+ HyenaDNA embedding (256) & LINSIGHT & worse & 0.12 & no difference \\
 & random columns (256) & worse & 0.04 & \textbf{worse} \\
 & shuffled embedding & worse & 0.29 & no difference \\
 & composition (4 numbers) & worse & 0.08 & no difference \\
+ autoencoder code (32) & LINSIGHT & better & 0.99 & no difference \\
 & random columns (32) & worse & 0.46 & no difference \\
 & shuffled code & better & 0.53 & no difference \\
+ autoencoder code (64) & LINSIGHT & better & 0.29 & no difference \\
+ autoencoder code (128) & LINSIGHT & better & 0.26 & no difference \\
+ positive control & LINSIGHT & worse & 0.25 & no difference \\
\midrule
Annotation code replacing the 48 annotations & LINSIGHT & better & $<0.0001$ & \textbf{better} \\
48 annotations + annotation code & LINSIGHT & better & 0.27 & no difference \\
\bottomrule
\end{tabular}
\end{table}

\subsection*{The withheld-genome test is too weak at this sample size}
To learn how large an effect the test can detect, we simulated evolutionary data on the real sites in
which the embedding truly shifts each site's selection score, and ran the full fit and test on each
simulated genome (Methods). The planted effect is expressed relative to the spread (standard deviation)
of LINSIGHT's own score across sites.

The test missed almost every planted effect (Table~\ref{tab:power}). Even when the embedding shifted
the true score by 23\% of the spread of LINSIGHT's score, it was detected in only one of five simulated
genomes. In simulated genomes with no effect, the test wrongly called the embedding better in 2 of 20
and random numbers better in 4 of 20, more often than the 1 in 40 expected by chance. At 0.1\% of
sites, an unchanged fit to withheld genome therefore does not show that sequence carries no
evolutionary information.

\begin{table}[t]
\centering\small
\caption{\textbf{Power of the withheld-genome test} in simulated genomes with a known sequence effect.}
\label{tab:power}
\begin{tabular}{@{}lcccccc@{}}
\toprule
Planted effect, \% of the spread of LINSIGHT's score & 0 & 1.5 & 2.9 & 5.8 & 12 & 23 \\
Simulated genomes & 20 & 5 & 5 & 5 & 5 & 5 \\
Embedding called better & 2 & 0 & 1 & 0 & 1 & 1 \\
\bottomrule
\end{tabular}
\end{table}

%==================================================================================================
\section*{Discussion}
At the 0.1\% sample, a HyenaDNA reading of the surrounding DNA did not improve LINSIGHT on either test.
The autoencoder code limited the damage done by the full embedding, but no more than random numbers of
its size did, so the compression appears to help by adding fewer columns, not sequence information. Only
the matched controls make this distinction possible. The withheld-genome test could not have detected a
modest effect, so the evolutionary question remains open. Alignment-aware language models, which see evolution directly,
are a natural next encoder.

\emph{Limitations.} One language model, one window size and one way of averaging its output;
polymorphism from 54 individuals; annotations on the hg19 genome and sequence on
hg38, joined by coordinate conversion; 1,564 disease variants; the HGMD read covers only models fitted
with the authors' procedure.

\paragraph{Remaining work.}
(i) Repeat every step on a 1\% sample of sites (which contains the 0.1\% sample), after widening the
range of the annotation prior, whose chosen value lies at the edge of the range tried, and measure
power again. (ii) A second, single HGMD read of the exactly fitted models; its plan is written and has
not been run. (iii) Confirm or reject the annotation-autoencoder result on the larger sample. (iv) Try
the alignment-aware encoder GPN-MSA, which needs a GPU. (v) Fill in the title,
contributions, funding and code link.

%==================================================================================================
\section*{Methods}

\subsection*{Method overview}
DeepINSIGHT asks whether the DNA sequence around a site tells LINSIGHT something its 48 annotations do
not. LINSIGHT scores each non-coding site with $\rho$, the probability that a new mutation there would be
removed by purifying selection. It computes $\rho$ from a linear combination of 48 genomic annotations and
learns the weights from evolutionary data alone: how often the site is polymorphic in humans and whether it
changed since the human--chimpanzee ancestor, compared with nearby neutral DNA. We keep that model and its
likelihood unchanged and add one block of sequence features to its input. A
pretrained genomic language model, HyenaDNA, reads the 1,001 bases around each site and summarises them as
256 numbers (the embedding). The embedding is appended to the 48 annotations either as it is or after a
sparse autoencoder has compressed it to 32, 64 or 128 numbers. Every model is fitted on the same sampled
sites with the same likelihood, and every added block is paired with a control block of the same width
that carries no sequence information (random or shuffled values). The fitted models then score two kinds
of held-out data: regulatory disease variants from HGMD against common variants (AUROC), and the
evolutionary data of genome blocks never used in fitting (likelihood). A sequence input helps only if it
beats both LINSIGHT and its own control. The pipeline is a chain of numbered steps
(\file{s00}--\file{s11}); each step reads named files, writes named files and a \file{checks.json} with
measured checks, and the next step refuses to start unless those checks pass.

\paragraph{Where files are.} Paths are relative to the project repository. Original inputs are given at
their location under \file{data/} (see \file{data/README.md} for how to obtain them); every file the pipeline writes is under
$R=$\,\file{outputs/rate_0.001/}, in a folder named after the step that wrote it
(e.g.\ \file{R/s02_sites/sites.npz}). Positions are on hg19 unless stated otherwise. A complete
file-by-file map is in \file{docs/DATA_FLOW.md}.

%--------------------------------------------------------------------------------------------------
\subsection*{Data Sources}

\paragraph{Description of the datasets used.} Nine original inputs are read and never modified
(Table~\ref{tab:sources}). Each is described below by what it is, where it came from, what one record looks
like, and what the pipeline uses it for.

\begin{table}[t]
\centering\footnotesize
\caption{\textbf{Original inputs.} ``Used by'' names the pipeline steps that read the file
(see Data Preprocessing). All nine are read-only.}
\label{tab:sources}
\begin{tabularx}{\textwidth}{@{}>{\raggedright\arraybackslash}p{2.55cm} >{\raggedright\arraybackslash}X >{\raggedright\arraybackslash}p{3.6cm} >{\raggedright\arraybackslash}p{3.0cm}@{}}
\toprule
Dataset & File in the repository & Obtained from & Used by \\
\midrule
INSIGHT database & \file{data/linsight/INSIGHT_database.ins} (text, 58.2\,GB) & LINSIGHT authors, \file{compgen.cshl.edu/LINSIGHT/INSIGHT_database.ins.gz} & \file{s02}: the sites and their evolutionary data \\
48 LINSIGHT annotations & \file{data/linsight/genomic_feature.bed} (BED, 28.2\,GB) & LINSIGHT authors, \file{compgen.cshl.edu/LINSIGHT/genomic_feature.bed.gz} & \file{s04}: the 48 input columns \\
Published LINSIGHT scores & \file{data/linsight/LINSIGHT.bw} (bigWig, 1.4\,GB) & LINSIGHT authors, \file{compgen.cshl.edu/LINSIGHT/LINSIGHT.bw} & \file{s01}, \file{s04}: reference scores \\
LINSIGHT programs & \file{third_party/linsight_cpp/} (\file{L-INSIGHT-prep.cpp}, \file{L-INSIGHT-fit.cpp}) & \file{github.com/CshlSiepelLab/LINSIGHT} (commit f28b163) & specification of parsing and fitting; baseline check \\
HGMD benchmark & \file{data/benchmark/} \file{hgmd_reg.pandas}, \file{unmatched.pandas}, \file{matched_by_tss.pandas}, \file{matched_by_region.pandas} & GWAVA v1.0 release, as used by LINSIGHT & \file{s01}: disease and common variants \\
hg19 genome & \file{data/genomes/hg19.fa} (FASTA, 3.2\,GB) & UCSC & \file{s01}--\file{s03}, \file{s06b}: base checks, composition \\
hg38 genome & \file{data/genomes/hg38.ml.fa} (FASTA, 3.1\,GB) & HyenaDNA data instructions & \file{s03}, \file{s06}: DNA windows \\
hg19$\to$hg38 chain & \file{data/genomes/hg19ToHg38.over.chain.gz} & UCSC liftOver & \file{s03}: coordinate conversion \\
HyenaDNA medium-160k & \file{data/hyenadna/hyenadna-medium-160k-seqlen/} (\file{weights.ckpt}, \file{config.json}); code in \file{third_party/hyena_dna/} & \file{huggingface.co/LongSafari/hyenadna-medium-160k-seqlen}; \file{github.com/HazyResearch/hyena-dna} (commit d553021) & \file{s05}, \file{s06}: the embedding \\
\bottomrule
\end{tabularx}
\end{table}

\emph{INSIGHT database.} This is the input file of the LINSIGHT paper: the
polymorphism and divergence data its authors prepared for INSIGHT and fitted
LINSIGHT to. It is a plain-text file with three kinds of line:
\begin{center}\footnotesize
\begin{tabular}{@{}l@{}}
\texttt{samples\ \ 108}\\
\texttt{block\ \ chr1:832500-837500\ \ theta\ \ 1.460239175e-03\ \ lambda\ \ 2.809890080e-03}\\
\texttt{site\ \ \ chr1:834081\ \ M\ \ 9.9991399050e-1}\\
\texttt{site\ \ \ chr1:834100\ \ L\ \ 9.9991399050e-1\ \ 7.6931202784e-5}
\end{tabular}
\end{center}
The first line gives the number of sampled chromosomes, 108 (54 unrelated individuals sequenced by Complete
Genomics). A \texttt{block} line opens a genomic interval and gives its neutral
polymorphism rate $\theta$ and neutral divergence rate $\lambda$, estimated by the authors from nearby
neutral DNA. Each \texttt{site} line that follows is one usable nucleotide in that interval: its position,
its polymorphism class in the 54 individuals (\texttt{M} monomorphic, \texttt{L} low-frequency or \texttt{H}
high-frequency polymorphism), and the probability that the ancestral allele, inferred from the chimpanzee,
orangutan and rhesus macaque genomes, is the major allele (and, for polymorphic sites, the minor allele). The
authors had already removed sites failing their quality filters (for example simple repeats, recent
transposable elements and segmental duplications); we applied no further filter
of this kind. The file contains 1,540,791,804 autosomal sites: 99.51\% \texttt{M}, 0.39\% \texttt{L} and
0.10\% \texttt{H}.

\emph{48 LINSIGHT annotations.} The genomic features LINSIGHT was trained on, as released by its authors:
conservation scores (phyloP, phastCons, SiPhy, CEGA), predicted transcription-factor and RNA binding sites,
and regional annotations from ENCODE, Roadmap Epigenomics, FANTOM5 and GENCODE.
The file is a tab-separated BED file sorted by chromosome and start. After a header line, each line is one
hg19 interval over which all 48 values are constant: \texttt{chrom}, \texttt{start} (0-based),
\texttt{end}, then 48 numbers, e.g.\ \texttt{chr1\ \ 9996\ \ 10003\ \ 0\ \ 0\ \ \dots\ \ 0.21\ \ 1\ \ 0\ \dots}
A position takes the values of the interval with $\texttt{start}\le p<\texttt{end}$; a position in no
interval has no annotations.

\emph{Published LINSIGHT scores.} The genome-wide LINSIGHT track ($\rho$ at every scored hg19 base) in
bigWig format, as distributed by the authors. It is not an input to any model; it is the reference that
our benchmark preparation and our recovered weights must reproduce.

\emph{LINSIGHT programs.} The authors' C\texttt{++} code: \file{LINSIGHT-prep} converts the INSIGHT database
to binary per-chromosome files, \file{LINSIGHT-fit} fits the weights. We read it as the specification
(parsing rules, constants, matching of sites to annotations, the fitting procedure) and ran it once on all
sites as the baseline check described under Evaluation.

\emph{HGMD benchmark.} The disease-variant benchmark of GWAVA, which LINSIGHT's
authors took directly from GWAVA for their own evaluation. It consists of four
Python-2 pandas tables (``pickles''); each row is one variant, indexed by its identifier (e.g.\
\texttt{CR000229} for an HGMD entry), with \texttt{chr}, \texttt{start} (0-based hg19), \texttt{end}
(${}=\texttt{start}+1$), 147 GWAVA annotation columns that we do not use, four one-hot columns
\texttt{seq\_A}, \texttt{seq\_C}, \texttt{seq\_G}, \texttt{seq\_T} giving the reference base, and the label
\texttt{cls}. \file{hgmd_reg.pandas} holds 1,614 regulatory disease variants from HGMD (\texttt{cls}${}=1$). The three negative sets are common variants (minor-allele
frequency above 1\% in the 1000 Genomes Project; \texttt{cls}${}=0$): \file{unmatched.pandas}
(161,400 variants drawn at random, 100 per disease variant), \file{matched_by_tss.pandas} (16,140 variants
matched to the disease variants by distance to the nearest transcription start site) and
\file{matched_by_region.pandas} (5,027 variants within 1\,kb of a disease variant). The three sets are
increasingly hard: the closer a negative lies to a disease variant, the more of its local context it shares.
The same folder holds two ClinVar sets, which we did not use (see below).

\emph{Genomes and chain file.} The hg19 FASTA (UCSC, soft-masked: repeats in lower case) is the assembly of
every input above. The hg38 FASTA is the one HyenaDNA was pretrained on, taken from HyenaDNA's data
instructions; it holds chr1--22 and chrX. The UCSC chain file \file{hg19ToHg38.over.chain.gz} maps hg19
coordinates to hg38.

\emph{HyenaDNA.} The released medium-160k checkpoint (\file{weights.ckpt}, 109.6\,MB; \file{config.json}: 8
layers, 256 numbers per base) and the official model code, copied unchanged with its BSD-3 licence.

\paragraph{Data origin and selection criteria.} Every input is a public release, used as distributed.
Each was compared with its upstream copy before use (step \file{s00}, \file{R/s00_inputs/checks.json}): the
INSIGHT database by the SHA-256 of the decompressed upstream file; the annotation file byte for byte with a
fresh download; the score track by its byte size (1,395,672,914); the chain file by its UCSC MD5 checksum;
the HyenaDNA weights by the SHA-256 recorded on HuggingFace; the two genomes by the length of every autosome
against UCSC. The inputs were chosen for these reasons. (i) The INSIGHT database, annotations, scores and
programs are exactly what LINSIGHT was built from, so any difference between our models and LINSIGHT is
due to the added sequence input and not to different data. (ii) The HGMD benchmark is the one LINSIGHT was
published on, with published AUROCs (0.897, 0.759 and 0.660 for the three negative sets) that our
preparation must reproduce. Only regulatory HGMD variants are positives, and only the three negative sets
LINSIGHT used are negatives. The ClinVar files in the same folder (194 disease variants) were not used; this
study evaluates on HGMD only. (iii) hg38 is used only for HyenaDNA, because the model was trained on hg38;
everything else stays on hg19, the assembly of all LINSIGHT inputs. (iv) HyenaDNA medium-160k is a released checkpoint pretrained on hg38 with single-base tokens, so every
base of the window is its own token; we used 1,001 of its 160,000 bases of context. (v) Sites are autosomal only, as the INSIGHT database is. The benchmark is kept as published,
including its 6,427 variants on chrX and chrY; the chrY variants drop out at the hg38 conversion because
the hg38 file has no chrY.

%--------------------------------------------------------------------------------------------------
\subsection*{Data Preprocessing}

Preprocessing runs as steps \file{s01}--\file{s06b} of the pipeline. Table~\ref{tab:funnel} shows how many
rows survive each step; Table~\ref{tab:files} lists every file the steps write, with its shape.

\begin{table}[t]
\centering\footnotesize
\caption{\textbf{Rows kept at each preprocessing step.} Sites are rows of \file{R/s02_sites/sites.npz};
benchmark variants are rows of \file{R/s01_benchmark/panel.csv.gz}. Rows removed after \file{s02} stay in
the files and are flagged, so every file keeps the same rows in the same order.}
\label{tab:funnel}
\begin{tabularx}{\textwidth}{@{}l >{\raggedright\arraybackslash}X r r@{}}
\toprule
Step & Rule & Sites & Benchmark variants \\
\midrule
input & all autosomal sites of the INSIGHT database; all rows of the four benchmark tables & 1,540,791,804 & 184,181 \\
\file{s02} sampling & keep a site if a fixed hash of (chromosome, position) is below 0.001 & 1,542,881 & -- \\
\file{s02} exclusion & remove sites within 1,000 bases of any benchmark variant & 1,362,112 & -- \\
\file{s03} hg38 conversion & keep unique, same-chromosome, same-strand conversions & 1,361,122 & 183,679 \\
\file{s04} annotations & keep rows inside an annotated interval & 1,342,159 & 181,398 \\
\file{s06}, \file{s06b} & keep rows with a finite embedding and composition (no row lost) & 1,342,159 & 181,398 \\
\file{s01} label conflict & remove the common variant at the position of a disease variant & -- & 181,397 \\
\midrule
usable rows & sites: 660,061 training, 69,412 \texttt{select}, 612,686 \texttt{eval} & 1,342,159 & 181,397 \\
\bottomrule
\end{tabularx}
\end{table}

\paragraph{Data cleaning and filtering.}
\emph{Benchmark (\file{s01_benchmark.py}).} The four pickles need pandas 0.13 under Python 2, so a small
script run in that environment exports only the columns we use (identifier, \texttt{chr}, \texttt{start},
\texttt{end}, \texttt{cls}, the four one-hot base columns) as text; the 147 GWAVA annotation columns are
dropped. The reference base is decoded from the one-hot columns (\texttt{N} if they are not exactly one 1).
The four tables are stacked into one panel, \texttt{label} = 1 for \file{hgmd_reg} and 0 otherwise, keeping
\texttt{start} as the hg19 position \texttt{pos19}. One common variant lies at the same position as a
disease variant; it is flagged (\texttt{label\_conflict}) and excluded from every evaluation. The published
LINSIGHT score is read from the bigWig at \texttt{pos19}. \emph{Checks:} the counts equal 1,614, 161,400,
16,140 and 5,027; the decoded reference base equals the hg19 base at \texttt{pos19} for 100\% of variants
(24.6\% one base off, so the coordinate convention is confirmed); the published scores reproduce the
published AUROCs to within 0.01 (0.898, 0.762, 0.664).

\emph{Sites (\file{s02_sites.py}).} The 58-GB INSIGHT database is parsed in 200 parallel chunks with the
rules of the authors' \file{LINSIGHT-prep}. Each site becomes five numbers: an allele-frequency code (0 for
\texttt{M}, 0.05 for \texttt{L}, 0.5 for \texttt{H}), the two ancestral-allele probabilities (the minor one
is 0 for \texttt{M}), $\theta a$ and $\lambda$ of its block, where $a=\sum_{k=1}^{107}1/k=5.2547$ is
Watterson's constant for 108 chromosomes and both rates are raised to at least $10^{-5}$, as the authors'
program does. \emph{Sampling:} a site is kept if a SplitMix64 hash of (chromosome index, position), read as
a number in $[0,1)$, is below the sampling rate 0.001. The hash does not depend on the order of reading, so
any run gives the same 1,542,881 sites, and the sample at a rate of 0.001 is contained in the sample at
0.01. \emph{Exclusion:} sites within 1,000 bases of any benchmark variant are removed (11.7\%), so that no
site's 1,001-base DNA window overlaps a benchmark variant's window and the benchmark cannot leak into
fitting; 1,362,112 sites remain. \emph{Split:} the genome is cut into 1-Mb hg19 blocks and whole blocks are
assigned at random, because neighbouring sites share their neutral rates and annotations. For the
likelihood analysis about half of the blocks are held out; one tenth of the held-out blocks form the
\texttt{select} set, used only to choose settings, and the rest the \texttt{eval} set, used only for the
final test (670,261 training, 70,332 \texttt{select}, 621,519 \texttt{eval} sites). The HGMD analysis used
an earlier split of the same sites with 10\% of blocks held out (1,212,074 training sites).
\emph{Checks:} the kept fraction is within 0.2\% of the rate (ratio 1.0014); the proportions of
\texttt{M}, \texttt{L} and \texttt{H} equal those of the whole file; the INSIGHT positions are 0-based (read
that way, 0.0\% of sites lie in a CpG dinucleotide, against 0.73\% and 0.79\% when read one base to either
side); no block is on both
sides of the split; every chromosome has sites in every set; the nearest kept site is 1,006 bases from any
benchmark variant.

\emph{Conversion to hg38 (\file{s03_liftover.py}).} Every site and benchmark variant is converted with the
UCSC chain file (pyliftover). A conversion is accepted only if it is unique, lands on the same chromosome and
keeps the strand. Rejected: 895 sites and 320 variants that map to the minus strand, 76 and 57 that map to
another chromosome, 19 and 123 that do not map, and 2 variants on chrY, which the hg38 file lacks.
\emph{Checks:} the base at the converted position equals the hg19 base for 100\% of sites and 99.85\% of
variants, and the whole 1,001-base windows agree at 99.98\% and 99.92\% of bases; the same test one base
off fails, as it must.

\emph{The 48 annotations (\file{s04_features.py}).} The 28-GB annotation file is read once, chromosome by
chromosome, and the 48 values are attached to every site and variant by the matching rule of the authors'
\file{LINSIGHT-fit} ($\texttt{start}\le p<\texttt{end}$). Sites are matched at their INSIGHT position, as
the authors' program does, and benchmark variants at \texttt{pos19}. Values were found for 98.6\% of sites
and 98.8\% of variants; the others lie outside every annotated interval and are dropped.
\emph{Check:} the published score is an exact function of the 48 values, so inverting its link
($u=-\log(-\log\rho/3)$) and regressing $u$ on the 48 values over the benchmark variants recovers the
published weights without any label ($R^2=1.000$). These weights reproduce the published score to within
0.002 at every variant and site, against about 20\% of rows when other rows' annotations are used.

\paragraph{Data formatting and preparation.}
\emph{DNA windows and tokens (\file{s06_embed.py}).} For each converted site and variant, the 1,001 hg38
bases centred on it (index 500 is the site) are read from \file{hg38.ml.fa}. Bases are mapped to HyenaDNA's
token ids (\texttt{A}, \texttt{C}, \texttt{G}, \texttt{T}, \texttt{N} $\to$ 7--11, case ignored, any other
character $\to$ 6) without start or end tokens, as in HyenaDNA's hg38 pretraining. How the window is turned
into the embedding is described under Computational Framework. The embeddings are stored with one row per
site (\file{R/s06_embed/sites_mean.npy}, $1{,}362{,}112\times256$, float32, 1.39\,GB) and per variant
(\file{R/s06_embed/panel_mean.npy}, $184{,}181\times256$); rows that did not convert are empty (NaN). A
fingerprint of all windows, the weights and the window length is stored with them, and the embeddings are
recomputed whenever it changes.

\emph{Composition control (\file{s06b_composition.py}).} Four numbers per 1,001-base hg19 window, centred on
\texttt{pos19}: GC fraction among \texttt{A}/\texttt{C}/\texttt{G}/\texttt{T}; CpG observed over expected,
$n_{CG}\,L/(n_C\,n_G)$ with $L=1{,}001$ (0 if $n_C$ or $n_G$ is 0); share of lower-case (repeat-masked)
bases; longest run of one base divided by $L$. They carry sequence information that needs no language
model.

\emph{Model inputs (\file{pipeline/arms.py}, checked by \file{s08_design.py}).} One function builds the input
matrix of every model, so all models see the same rows in the same order: the 48 annotations followed by at
most one added block (Table~\ref{tab:inputs}). Each added column is standardised with the mean and standard
deviation of the training sites, and the same values are applied to \texttt{select} and \texttt{eval} sites
and to benchmark variants. A row is used only if it has annotations, an accepted hg38 conversion, an
embedding and composition values (Table~\ref{tab:funnel}); every model uses the same usable rows.
\emph{Check:} for 300 random training sites and 300 benchmark variants, every model's input row was rebuilt
from the original files with independent code (a separate annotation lookup, a fresh conversion, a fresh
HyenaDNA run, reloaded autoencoders) and matched the stored row.

\begin{table}[t]
\centering\footnotesize
\caption{\textbf{Files written by preprocessing.} $n_s=1{,}362{,}112$ sites (rows in the order of
\file{sites.npz}); $n_p=184{,}181$ benchmark variants (rows in the order of \file{panel.csv.gz}). All paths
are under $R$.}
\label{tab:files}
\begin{tabularx}{\textwidth}{@{}>{\raggedright\arraybackslash}p{4.9cm} >{\raggedright\arraybackslash}X@{}}
\toprule
File & Content and shape \\
\midrule
\file{s01_benchmark/panel.csv.gz} & table, $n_p\times11$: \texttt{key, source, row\_id, chrom, pos19, end, cls, label, ref, label\_conflict, linsight\_published} \\
\file{s02_sites/sites.npz} & \texttt{chrom} ($n_s$), \texttt{pos\_ins} (INSIGHT position), \texttt{pos19}, \texttt{evol} ($n_s\times5$: frequency code, P(ancestral${}={}$major), P(ancestral${}={}$minor), $\theta a$, $\lambda$), \texttt{role} (train / select / eval) \\
\file{s03_liftover/sites_lift.npz}, \file{panel_lift.npz} & \texttt{pos38} and \texttt{status} (\texttt{ok}, \texttt{minus\_strand}, \texttt{other\_chrom}, \texttt{unmapped}, \texttt{no\_hg38\_chrom}) per row \\
\file{s04_features/sites_features.npz}, \file{panel_features.npz} & \texttt{F} ($n\times48$, float32) and \texttt{found} per row \\
\file{s04_features/published_weights.npz} & \texttt{rho\_coef}: intercept and 48 recovered published weights \\
\file{s06_embed/sites_mean.npy}, \file{panel_mean.npy} & embeddings, $n\times256$ float32 (NaN where not converted) \\
\file{s06b_composition/sites_comp.npy}, \file{panel_comp.npy} & composition, $n\times4$ float32 \\
\file{s07_sae/codes_<size>_s<seed>_sites.npy}, \file{..._panel.npy} & autoencoder codes, $n\times d$ ($d=32,64,128$) for seeds 0--4; models in \file{s07_sae/*.pt} \\
\bottomrule
\end{tabularx}
\end{table}

%--------------------------------------------------------------------------------------------------
\subsection*{Computational Framework}
\paragraph{Overview of the computational approach.} The method has three parts. (1) A frozen language
model turns the DNA around a site into a fixed-length vector. (2) An optional sparse autoencoder compresses
that vector. (3) LINSIGHT's statistical model uses the 48 annotations plus the added vector to predict the
site's evolutionary data. Only part (3) is trained against evolutionary data, and no part sees disease
labels.

\paragraph{Description of the model or methods used.}
\emph{LINSIGHT.} LINSIGHT extends INSIGHT, a
population-genetic model of one site. INSIGHT assumes that a site is either neutral or under purifying
selection. A neutral site becomes polymorphic in humans, and diverges from the human--chimpanzee ancestor,
at the local neutral rates $\theta$ and $\lambda$ of its block. A selected site cannot carry a common
polymorphism, and keeps rare variants and substitutions only at a reduced rate set by a second selection
term, $\gamma'$. Each site enters as the five numbers of \file{sites.npz}. The probability of the observed
data is
\begin{equation}
P(\text{site}) = \big(A\,(B + D\,\beta_0\,\gamma') - F\big)\,\rho + F ,
\label{eq:lik}
\end{equation}
where $F$ is the probability of the data if the site is neutral, $A(B + D\beta_0\gamma')$ the probability
if it is selected ($A=0$ for a common polymorphism), and $A,B,D,F$ are fixed functions of the five numbers
with constants $\beta=(0.757, 0.205, 0.038)$, computed exactly as in \file{L-INSIGHT-fit.cpp}. LINSIGHT
makes the two selection quantities functions of an input vector $\mathbf{x}$:
\begin{equation}
\rho(\mathbf{x})=\exp\!\big(-3\,e^{-(b+\mathbf{w}^{\top}\mathbf{x})}\big),\qquad
\gamma'(\mathbf{x})=\sigma(c+\mathbf{v}^{\top}\mathbf{x}),
\label{eq:linsight}
\end{equation}
with $\sigma$ the logistic function. The weights $(b,\mathbf{w},c,\mathbf{v})$ are found by maximising
$\sum\log P$ over training sites, so a model with $m$ input columns has $2(m+1)$ parameters (98 for the 48
annotations). In effect LINSIGHT is a generalised linear model whose labels are replaced by a mixture
likelihood: an annotation earns a positive weight if sites where it is high show fewer polymorphisms and
substitutions than their neutral neighbours. The output $\rho$ is the site's score.

\emph{HyenaDNA.} HyenaDNA is a genomic language model: a network trained on the
human reference genome (hg38) to predict each nucleotide from the nucleotides before it, with one token per
base. In place of attention, each layer uses the Hyena operator: very long
convolutions whose filters are generated by a small network over positions, alternated with element-wise
multiplications by projections of the input (gating). The convolutions are computed with the fast Fourier
transform, which makes contexts of up to 160,000 bases affordable. We used the released medium-160k model
(8 layers, a 256-number state per base) without further training. The model is causal: the state at a base
summarises that base and the bases before it. \emph{From HyenaDNA to an embedding:} each 1,001-token window
(see Data formatting) is passed through the frozen model in evaluation mode; the last layer gives a
$1{,}001\times256$ matrix, and its average over positions is the embedding $\mathbf{h}\in\mathbb{R}^{256}$.
We averaged rather than taking the state at the site because a causal state at the centre has seen only the
500 bases on its left; the average uses both flanks. \emph{Checks (\file{s05_embedder.py}):} the model
predicts the next base of real hg38 DNA in 43.5\% of cases against 25.3\% on random sequence; output is
identical across runs, batch sizes and processes; changing base 501 does not change the state at base 500.

\emph{Integration of LINSIGHT and HyenaDNA.} Each model's input is $\mathbf{x}=[\,\mathbf{a}\,;\,\mathbf{e}\,]$:
the 48 annotations $\mathbf{a}$ followed by at most one added block $\mathbf{e}$ (Table~\ref{tab:inputs}),
which is the embedding $\mathbf{h}$, a sparse code $\mathbf{z}$ of it, or a control. No model uses both
$\mathbf{h}$ and $\mathbf{z}$. In equation~\ref{eq:linsight} the added block contributes extra weights,
$b+\mathbf{w}_a^{\top}\mathbf{a}+\mathbf{w}_e^{\top}\mathbf{e}$, so the sequence can only shift a site's
probability of selection up or down; with $\mathbf{w}_e=\mathbf{0}$ the model is exactly LINSIGHT.
\emph{Sparse code (\file{s07_sae.py}).} Appending all 256 numbers raises the parameter count from 98 to 610
in a model fitted to a 0.1\% sample of sites. A TopK sparse autoencoder compresses the
embedding to a code in which only $k$ of $d$ entries may be nonzero:
$\mathbf{z}=\mathrm{TopK}_k\big(\mathrm{ReLU}(W_{\mathrm{enc}}(\tilde{\mathbf{h}}-\mathbf{b}_{\mathrm{dec}})+\mathbf{b}_{\mathrm{enc}})\big)$
and $\hat{\mathbf{h}}=W_{\mathrm{dec}}\mathbf{z}+\mathbf{b}_{\mathrm{dec}}$, where $\tilde{\mathbf{h}}$ is the
embedding centred on the training mean and divided by one scale factor, and the columns of
$W_{\mathrm{dec}}$ have unit length. Training minimises the squared rebuilding error plus the auxiliary term
of Gao et al. that revives code entries unused for 500 batches (weight 1/32). Three sizes,
$(d,k)=(32,8)$, $(64,16)$ and $(128,32)$, each with five random starts, were trained on the embeddings of
the 669,704 training sites only, then applied to every site and variant. A code of 32--128 entries adds
64--256 parameters, and its entries can be inspected one by one. \emph{Checks:}
every code entry is used, and every autoencoder rebuilds held-out embeddings better than principal
component analysis with as many components (for 32 entries, 4.9\% against 7.4\% of the variance left
unexplained). A second autoencoder of the same kind, $(32,8)$, was trained on the 48 annotations.

\begin{table}[t]
\centering\footnotesize
\caption{\textbf{Model inputs.} Each row is one model; ``Added block'' follows the 48 annotations. The code
name is the model's name in the pipeline files (e.g.\ \texttt{emb\_s0} in \file{params.npz}).}
\label{tab:inputs}
\begin{tabularx}{\textwidth}{@{}l l c X@{}}
\toprule
Model & Code name & Added block & What the added block contains \\
\midrule
LINSIGHT (48 annotations) & \texttt{repl} & none & -- \\
+ HyenaDNA embedding & \texttt{emb} & 256 & \file{sites_mean.npy} / \file{panel_mean.npy} \\
+ autoencoder code & \texttt{sae32k8}, \texttt{sae64k16}, \texttt{sae128k32} & 32, 64, 128 & the code of the embedding, used instead of it \\
+ random columns (control) & \texttt{noise\_emb}, \texttt{noise32/64/128} & same width & $\mathcal{N}(0,1)$ numbers, fixed seeds, independent of the site \\
+ shuffled embedding or code (control) & \texttt{shuf\_emb}, \texttt{shuf\_sae*} & same width & real values given to the wrong sites \\
+ composition (control) & \texttt{comp} & 4 & \file{sites_comp.npy} / \file{panel_comp.npy} \\
+ positive control & \texttt{posctl} & 1 & the annotation with the largest published weight, moved into the added block \\
annotation autoencoder & \texttt{fsae\_codes}, \texttt{fsae\_plus} & 32 & a code of the 48 annotations, replacing them or added to them \\
\bottomrule
\end{tabularx}
\end{table}

\paragraph{How the model or computational method was applied.} Every model was fitted in two ways.
(i) The \emph{authors' procedure} (\file{s09_train.py}), an exact re-implementation of \file{LINSIGHT-fit}:
all parameters start at zero, sites are visited in random batches of 100, AdaGrad
adapts the step size of each parameter (learning rate $10^{-3}$ for the $\rho$ weights and $10^{-2}$ for
the $\gamma'$ weights), and the parameters after the 20th pass are kept; $\mathbf{x}$ enters both $\rho$
and $\gamma'$. This is the procedure behind the published track and was used for the HGMD analysis: three
runs per model with different site orders on the 10\% split, scoring each variant by the average $\rho$ of
the three runs. (ii) The \emph{exact fit} (\file{s09b_mle.py}), used for the withheld-genome test:
penalised maximum likelihood run to convergence with L-BFGS on the 51\% split. The
penalty is a Gaussian prior on each weight; its width for the 48 annotation weights was chosen once, for
LINSIGHT, on \texttt{select} sites and kept for every model, and its width for the added weights was chosen
per model on \texttt{select} sites. Added columns enter $\rho$ only. Each fit was started from zero and from
LINSIGHT's solution, keeping the better; models with a random block were fitted five times. The
autoencoder size was chosen on \texttt{select} sites, never with benchmark labels: 64 for the HGMD analysis
and 32 for the exact fit. Fitted weights are in \file{R/s09_train/params.npz} and
\file{R/s09b_mle/params.npz} (one array per model and seed). \emph{Check:} every exact fit converged, and no
model that contains LINSIGHT fits the training data worse than LINSIGHT. To score a benchmark variant, the
fitted weights are applied to that variant's own $\mathbf{x}$.

\paragraph{Key implementation choices and parameters.} Table~\ref{tab:settings} lists every setting.

\begin{table}[t]
\centering\footnotesize
\caption{\textbf{Settings.} Values chosen on \texttt{select} sites are marked; nothing was tuned on the benchmark.}
\label{tab:settings}
\begin{tabularx}{\textwidth}{@{}l X@{}}
\toprule
Component & Setting \\
\midrule
Sites & autosomal INSIGHT sites sampled at 0.1\% by a hash of position; sites within 1 kb of a benchmark variant removed \\
Splits & whole 1-Mb blocks; HGMD analysis 90/10 train/held-out; likelihood analysis 49/51, of which 10\% \texttt{select} \\
Sequence window & 1,001 bp of hg38 centred on the site; tokens A, C, G, T, N; no special tokens \\
HyenaDNA & medium-160k, 8 layers, 256-d state, weights frozen; last layer averaged over the window \\
Autoencoder & TopK; $(d,k)\in\{(32,8),(64,16),(128,32)\}$; Adam, lr $10^{-3}$, batch 4,096, 10 passes, 5 seeds; size chosen on \texttt{select} (64 for HGMD, 32 for exact fit) \\
Authors' procedure & zero start; batches of 100; AdaGrad, lr $10^{-3}$ ($\rho$) and $10^{-2}$ ($\gamma'$); 20 passes; 3 runs per model \\
Exact fit & L-BFGS to convergence; Gaussian prior widths: annotations $\{0.03,\dots,10\}$ chosen once on \texttt{select}; added block $\{0.003,\dots,10\}$ chosen per model on \texttt{select} \\
Controls & random $\mathcal{N}(0,1)$ columns, shuffled embedding or code, 4 composition values; same width as the block they control \\
Statistics & paired cluster bootstrap, 2,000 resamples; 100-kb regions (HGMD), 1-Mb blocks (withheld genome) \\
\bottomrule
\end{tabularx}
\end{table}

%--------------------------------------------------------------------------------------------------
\subsection*{Computational Reproducibility}
\paragraph{Software and computing environment.} Python 3.10.20 with PyTorch 2.12.0 (CPU build), NumPy 2.2.6, SciPy 1.15.3,
scikit-learn 1.7.2, pandas 2.3.3, pyliftover, pyBigWig and Matplotlib (conda environment \file{hyena-env}); the benchmark pickles are read with Python
2.7.13 and pandas 0.13.1 (environment \file{legacy_pandas}). HyenaDNA's official model code and weights;
the authors' LINSIGHT programs as reference. Everything ran on one 112-core CPU server without a GPU; the
HyenaDNA embedding was the slowest step (100 processes, 16 windows per batch).

\paragraph{Code and reproducibility details.} The pipeline is in this repository, with the steps in
\file{pipeline/steps/}, the shared code in \file{src/lctx/} and the paths of the original inputs in
\file{src/lctx/config.py}. The script \file{pipeline/run_all.sh}, called with the sampling rate (0.001), runs every step in
order and stops at the first failed check; \file{pipeline/run_step.py} runs a single step. All
randomness is seeded, inputs are verified against their upstream checksums, and each step writes its checks
with measured values to \file{R/<step>/checks.json}. Our likelihood code reproduces \file{LINSIGHT-fit}
parameter for parameter, and LINSIGHT re-fitted with the authors' program on all sites reproduces the
published AUROCs. Every check passed except two in the power simulation, which are the low power reported in
the Results.

%--------------------------------------------------------------------------------------------------
\subsection*{Evaluation and Benchmarking}
\paragraph{How performance was assessed.} \emph{HGMD benchmark (\file{s10_evaluate.py}).} For each of the
three negative sets, the AUROC was computed on all usable disease variants (1,564) and that set's usable
common variants (159,765 unmatched, 15,346 TSS-matched, 4,722 region-matched): the probability that a
randomly chosen disease variant has a higher $\rho$ than a randomly chosen common one. The
analysis plan (models, comparisons, number of runs, bootstrap) was fixed and recorded before the labels were
opened (\file{R/s10_evaluate/lock.json}), and the code refuses a second read of the same plan; the benchmark
has been read once. Results are in \file{R/s10_evaluate/results.json}. \emph{Withheld genome
(\file{s09b_mle.py}).} For each \texttt{eval} site we computed $\log P$ of its observed data under each fitted
model; two models are compared by the mean difference over sites, after averaging the five fits of models
with a random block. A better model gives the real data a higher probability. \emph{Power
(\file{s09c_groundtruth.py}).} Using the real sites, a true score was set to LINSIGHT's fitted score plus an
effect size times a fixed random combination of the embedding's 256 numbers; new polymorphism and divergence
data were drawn from equation~\ref{eq:lik} with that score; every model was refitted exactly as on the real
data and the withheld-genome test applied: 20 simulated genomes without an effect and five at each of five
effect sizes.

\paragraph{Comparison methods or baselines.} The reference is LINSIGHT with the 48 annotations alone,
fitted with the same procedure, sites and split as every other model. As a check of the baseline itself,
the authors' programs were run on all sites (\file{LINSIGHT-prep} on \file{INSIGHT_database.ins}, then
\file{LINSIGHT-fit} with \file{genomic_feature.bed}, 20 passes) and their final weights applied to the
benchmark variants; the published score track was also scored directly. Each sequence model is compared
with LINSIGHT and with its control of the same width; the positive control checks that the added-block path
can carry a real signal.

%--------------------------------------------------------------------------------------------------
\subsection*{Statistical Analysis}
\paragraph{Statistical tests and summary measures.} The summary measures are the AUROC per model and
negative set, and the mean per-site difference in $\log P$ on withheld sites. Differences between two
models are paired: both are scored on the same variants or sites. Uncertainty comes from a cluster
bootstrap: whole regions are resampled with replacement 2,000 times (100-kb
regions for HGMD, 1-Mb blocks for the withheld genome, because nearby variants and sites are correlated),
the difference is recomputed each time, and the middle 95\% of the results is the 95\% interval. A
two-sided $p$-value is the probability of a difference at least this large if the two models were equally
good.

\paragraph{Criteria used to evaluate results.} A sequence input counts as an improvement only if its
difference from LINSIGHT and from its same-width control are both positive with 95\% intervals that
exclude zero. We did not correct for the 22 withheld-genome comparisons; about one $p$ below 0.05 is
expected among them by chance, so single uncorrected results are reported as such.

%==================================================================================================
{\small
\noindent\textbf{Author Contributions:} [we will fill this in soon]

\noindent\textbf{Funding:} [Add funding] The content is solely the responsibility of the authors and
does not necessarily represent the official views of the National Institutes of Health.

\noindent\textbf{Code Availability Statement:} The DeepINSIGHT code, with per-step checks, unit tests and
the data-flow document, can be found in the following repository: [link-here-eventually]. HGMD-derived benchmark files are not redistributed and
must be obtained from their original sources.

\noindent\textbf{Competing Interest Statement:} The authors declare no competing interest.\par}


\end{document}
